\section{Blind one-sided separation on the positive half-line}
\label{sec:riesz}

The article uses a positive prime-frequency weight for multiplicity differences of the input $L$-functions in \texttt{zer:thm:short}. Here the zeros belong to the coupled outputs. A single arithmetic trace supplies a lower bound for their total discrepancy. A nonnegative trigonometric polynomial combines several traces to force a linear discrepancy in each direction, even when one conductor product is much larger than the other.

\begin{theorem}\label{thm:directional}
Assume the critical-line input. If two couplings of the same size have different character multisets, there is $d>0$, depending on the pair, such that for all sufficiently large $T$,
\[
 \min\bigl\{\Delta_{(0,T]}(P_A,P_B),\Delta_{(0,T]}(P_B,P_A)\bigr\}\ge dT .
\]
\end{theorem}
\begin{proof}
For real $\omega$ put $J_T(\omega)=\sum_{0<\gamma\le T}(\nu_A(\gamma)-\nu_B(\gamma))e^{i\omega\gamma}$. For each fixed real $x>1$, the proof of Theorem~\ref{thm:landau} gives
\[
 J_T(\log x)=T b(x)+O_{A,B,\mathcal A,\mathcal B,x}(T/\log T),
 \quad
 b(x)=\begin{cases}
 -\dfrac{\log p}{2\pi\sqrt x}(S_{\mathcal A}(x)-S_{\mathcal B}(x)) &x=p^j,\quad j\ge1,\\[3pt]
 0 &\text{otherwise}.
 \end{cases}
\]
Indeed the Poisson argument applies to every $\log x>0$, and the distance from $\log x$ to the nearest different $\log l$, $l\ge2$ an integer, is positive for fixed $x$. The only possible resonant Euler term comes from $l=x$, with coefficient zero unless $x$ is a prime power. At the negative frequency $J_T(-\log x)=\overline{J_T(\log x)}$, so the main coefficient is $\overline{b(x)}$ and the same error estimate holds. The ordinary one-sided Riemann--von Mangoldt count for each primitive character is $T\log T/(2\pi)+O_\chi(T)$. Since the two lists have equal size, the coupled count comparison gives a constant $C_0$ with $|J_T(0)|\le C_0 T$ for $T\ge3$.

Choose a unit class $a$ modulo a common multiple of all conductors on which the two character sums differ. Such a class exists by character orthogonality. Along primes $p\equiv a$ the absolute value of that difference is a fixed positive number. The prime number theorem in this fixed progression gives $\sum_{p\equiv a}|b(p)|=\infty$; for example the primes in $[X,2X]$ already contribute $\gg\sqrt X$ for large $X$. Choose a finite set $\mathcal S$ of these primes with $\sum_{p\in\mathcal S}|b(p)|>C_0+2$, and choose $\theta_p$ with $e^{i\theta_p}b(p)=|b(p)|$. Consider the two Riesz products
\[
 W_\pm(t)=\prod_{p\in\mathcal S}\bigl(1\pm\cos(t\log p+\theta_p)\bigr),\qquad
 0\le W_\pm(t)\le 2^{|\mathcal S|}.
\]
Their constant terms are $1$ by unique factorization. Each term involving at least two distinct primes has frequency $\log r$, where $r$ is a positive rational with exponents $+1$ or $-1$ at two or more distinct primes. Neither $r$ nor $r^{-1}$ is a positive integral prime power. Consequently its trace has zero main term. The singleton terms contribute $\pm T|b(p)|$ for each prime, and the finite number of errors is $o(T)$. It follows that
\[
 \sum_{0<\gamma\le T}(\nu_A(\gamma)-\nu_B(\gamma))W_\pm(\gamma)
 =J_T(0)\ \pm T\sum_{p\in\mathcal S}|b(p)|+o(T).
\]
For large $T$ the plus expression is at least $T$ and the minus expression is at most $-T$. Nonnegativity and the upper bound on the weights give both claimed lower bounds, with $d=2^{-|\mathcal S|}$.
\end{proof}

\begin{corollary}[Positive ordinates, unrestricted surplus]\label{cor:optimal-heading}
Let the true list size be $n$. Suppose a competitor has at most $n$ vertices and the number of true positive-ordinate multiplicities missing from it through $T$ is $o(T/(\log T)^{n+1})$. Under the critical-line input, the character multisets and grouped determinant polynomials coincide. The competitor may have arbitrary conductor, matrix norm, and surplus multiplicities.
\end{corollary}
\begin{proof}
A smaller competitor is ruled out by the leading term of the ordinary positive-ordinate counts and the coupled count bound. At equal size Theorem~\ref{thm:directional} identifies the labels, then Theorem~\ref{thm:minortrace} identifies the grouped polynomial. A larger direct-sum competitor shows that the vertex restriction under one-sided observation is necessary.
\end{proof}
